\section{Appendix: Scope and Qualifications}

This concise study compares the two readings developed in \work{The Feast
and the New Person}, the expansive study of these same ten propers.
It concerns the Nineteenth Sunday after Pentecost, class II, in the
1962 Roman Missal, pp.~411--412, for the general-calendar occurrence
on 4 October 2026. Local calendars, a commemoration of St Francis and
the optional external solemnity of the Rosary are excluded. No present
permission to celebrate with this book is determined here. The complete
appointed texts belong to the expansive study; this companion gives
their inventory and exposition. The Creed and Trinity Preface are
appointments, not additional texts reproduced here.

The two whole-formulary readings are editorial syntheses of the
attributed exegesis; Schuster directly expounds the matching Mass.
The reception sweep is bounded, not exhaustive. Historical translations
and Latin transcriptions are not fresh critical editions; Aquinas's
inspected OCR supports paraphrase only. The historical dossier preserves
competing chronology and authorship judgments and the unresolved Gospel
event date. Its broad Psalter boundary dates no individual occasion.
Evidence, text collation and interpretive limits are recorded in
\texttt{propers/verified.md}, \texttt{research/scope.md} and
\texttt{research/interpretations.md}.

Biblical wording is Douay--Rheims--Challoner; Latin identifiers belong
to the collated 1962 text with its public-domain 1862 antecedent.
The 1962 edition itself is not declared public domain. No new prayer
or translation is supplied. The work is an AI-assisted Catholic study
aid, not an official liturgical edition or ecclesiastical judgment.

\section{References}
\begin{description}[style=nextline,leftmargin=0pt,labelsep=0pt,itemsep=.35em]
\item[Appointed texts and biblical context]
\work{Missale Romanum}, Vatican editio typica, 1962, pp.~411--412,
nn.~1679--1688; Pustet, Ratisbon, 1862, pp.~351--352, Latin antecedent.
Douay--Rheims--Challoner, Gutenberg eBook 1581, Psalms 77, 104, 118,
137, 140; Matthew 21--22; Ephesians 4--5.
\url{https://media.churchmusicassociation.org/pdf/missale62.pdf};
\url{https://www.gutenberg.org/ebooks/1581}.

\item[Patristic exegesis]
Augustine, Sermon 90.1--10 (NPNF I.6, English Sermon XL);
\work{Enarrationes in Psalmos} 77.1--3, 104.1--3, 118 opening
Aleph exposition 4--5, 137.10--13, 140.2--4 (NPNF I.8; English
Psalm headings 78, 105, 119, 138, 141).
\url{https://www.ccel.org/ccel/schaff/npnf106.html};
\url{https://www.ccel.org/ccel/schaff/npnf108.html}.
Gregory the Great, \work{Homiliae in Evangelia} II.38.1--3, 7,
9--14, historical Latin, Wikisource transcription:
\url{https://la.wikisource.org/wiki/Homiliarum_in_Evangelia}.
John Chrysostom, \work{Homilies on Ephesians} XIII, on 4:22--24,
and XIV, on 4:25--28, NPNF I.13:
\url{https://www.ccel.org/ccel/schaff/npnf113.html}.

\item[Later exegesis and liturgical reception]
Robert Bellarmine, \work{A Commentary on the Book of Psalms},
trans. John O'Sullivan, 1866, Psalms 77:1--8, 104:1--2,
118:4--8, 140:1--3; historical English, paraphrased:
\url{https://www.ecatholic2000.com/bellarmine/commentary.shtml}.
Thomas Aquinas, \work{Super Ephesios} 4, lectio 7, Dessain 1857,
vol.~II, p.~329, inspected Latin OCR, paraphrased.
Blessed Ildefonso Schuster, \work{The Sacramentary}, English 1927,
vol.~III, ``Nineteenth Sunday after Pentecost,'' pp.~171--174.

\item[Chronology and historical setting]
\work{Catholic Encyclopedia}: Paulin Ladeuze, ``Epistle to the
Ephesians'' (1909), date/place and destination;
\url{https://www.newadvent.org/cathen/05485a.htm}; Ferdinand Prat,
``St. Paul'' (1911), chronology table;
\url{https://www.newadvent.org/cathen/11567b.htm}; E. Jacquier,
``Gospel of St. Matthew'' (1911), date/place;
\url{https://www.newadvent.org/cathen/10057a.htm}; Alfred Durand,
``The Synoptics'' (1912), ``Origin,'' Aramaic original and Greek
rendering; \url{https://www.newadvent.org/cathen/14530a.htm}.
USCCB, NABRE introductions to Psalms and Matthew, and Ephesians
paragraphs 3--4, consulted 28 September 2026:
\url{https://bible.usccb.org/bible/psalms/0};
\url{https://bible.usccb.org/bible/matthew/0};
\url{https://bible.usccb.org/bible/ephesians/0}.
These introductions supply historical judgments, not liturgical English.
\end{description}
